\documentclass[../../../include/open-logic-section]{subfiles}

\begin{document}

\olfileid{sth}{spine}{zf}
\olsection{$\Z$ आणि $\ZF$: एक महत्त्वाचा टप्पा}

पायाभूततेमुळे आपण आणखी एक महत्त्वाचा टप्पा गाठतो. आपण $\Zminus$ आणि $\ZFminus$ या उपपत्ती पाहिल्या; त्यांतून एक विशिष्ट गोष्ट ``वगळलेली'' आहे असे म्हटले होते. ती गोष्ट म्हणजे पायाभूतता. म्हणून:

\begin{defn}
$\Zminus$ मध्ये पायाभूतता जोडली की $\Z$ ही उपपत्ती मिळते. म्हणून तिची स्वयंसिद्धके विस्तारात्मकता, संयोग, जोड्या, घातसंच, अनंतता, पायाभूतता आणि विभक्तीकरण-रूपबंधातील सर्व स्वयंसिद्धके ही आहेत.

$\ZFminus$ मध्ये पायाभूतता जोडली की $\ZF$ ही उपपत्ती मिळते. दुसऱ्या शब्दांत, $\Z$ मध्ये प्रतिस्थापनाची सर्व स्वयंसिद्धके जोडली की $\ZF$ मिळते.
\end{defn}

तरी एक प्रश्न पडू शकतो: $\ZFminus$ मध्ये नियमितता पायाभूततेशी समतुल्य आहे आणि नियमिततेचे समर्थन स्पष्ट आहे, तर थेट नियमिततेला मूलभूत स्वयंसिद्धक म्हणून न घेता पायाभूततेचा आडमार्ग का घ्यावा?

ऐतिहासिक कारणे (कोणी कोणते तत्त्व कधी मांडले, ही) बाजूला ठेवली, तर मुख्य कारण असे: $V_\alpha$ च्या व्याख्येचा वापर न करता पायाभूतता मांडता येते. पण त्या व्याख्येसाठी \olref[recursion]{sec} मधील सर्व काम आवश्यक होते: ती व्याख्या योग्य आहे हे दाखवण्यासाठी अनंतपार पुनरावर्तन सिद्ध करावे लागले. आणि त्या पुराव्यात \emph{प्रतिस्थापन} वापरले. म्हणून $\ZFminus$ च्या संदर्भात पायाभूतता आणि नियमितता समतुल्य असल्या, तरी $\Zminus$ च्या संदर्भात त्या समतुल्य नाहीत.

खरे तर प्रश्न या साध्या निरीक्षणापेक्षा अधिक गंभीर आहे. या पुस्तकाच्या व्याप्तीपलीकडील निष्कर्ष असा की $\Zminus$ आणि $\Z$ या दोन्ही उपपत्ती $V_\alpha$ परिभाषित करण्यास अपुऱ्या आहेत. म्हणून फक्त $\Z$ मध्ये काम करत असाल, तर आपण मांडलेल्या रूपातील नियमिततेला \emph{अर्थच} नाही. म्हणून नियमिततेऐवजी पायाभूतता हेच आपले अधिकृत स्वयंसिद्धक आहे.

यापुढे अन्यथा सांगितले नसेल, तर आपण अधिक काही न सांगता $\ZF$ मध्ये काम करू.

\end{document}
